\appendix
\chapter{ملاحق الفوائد المستنبطة}

هذا القسم يجمع الفوائد والدرر المتناثرة التي تم استنباطها وتقييدها خلال دراسة الكتاب، وقد تم تصنيفها موضوعياً لتكون دليلاً سريعاً وكشافاً لطالب العلم، وتزداد حصيلتها تباعاً مع التقدم في دروس التفسير.

\section{الفوائد التفسيرية وعلوم القرآن}
\begin{itemize}
    \item \textbf{أهمية سورتي الفاتحة والبقرة لاكتشاف منهج الطبري:} تفسير الطبري لهاتين السورتين هو أهم ما يُستكشف به منهج المؤلف، فقد توسع فيهما توسعاً كبيراً جداً، وكان يعزو إليهما لاحقاً. فمن درسهما كأنه استكشف أهم ما عند الطبري (اللقاء الثاني).
    \item \textbf{ضابط المكي والمدني في النداء:} الآية المفتتحة بـ (يا أيها الناس) ليست حصرية في السور المكية كما ظن بعض المتأخرين، فسورة البقرة مدنية وفيها هذا النداء (اللقاء التاسع).
    \item \textbf{تقسيم علوم القرآن:} تنقسم إلى علوم تضمنها النص نفسه (كالعام والخاص والمحكم والمتشابه)، وعلوم محتفة بالنص وتأخذ صبغة تاريخية (كالمكي والمدني وأسباب النزول). والطبري في خطبة كتابه عَنَى النوع الأول (اللقاء الثاني).
    \item \textbf{مفهوم الاستيعاب عند الطبري:} عندما يصف الطبري كتابه بأنه جامع، فهو يعني استيعاب الأقوال والمذاهب ومواضع الإجماع والاختلاف، وليس بالضرورة استيعاب كل المرويات والأسانيد المكررة في ذات المعنى (اللقاء الثاني).
    \item \textbf{منهجية توجيه الأقوال:} لا يكتفي الطبري بذكر الأقوال، بل يُعنى بـ (توجيهها) بذكر العلة والدليل والأصل الذي بُني عليه القول ولو كان ضعيفاً ليرده (اللقاء الثاني).
    \item \textbf{الفرق بين التفاسير المتقدمة وتفسير الطبري:} اعتمد من قبله (كمقاتل بن سليمان ويحيى بن سلام) إما على الاختيار أو الجمع المجرد، أما الطبري فقد أبدع الطريقة النقدية التحليلية المستندة لقواعد الترجيح (اللقاء الثاني).
    \item \textbf{وجه الإعجاز في القرآن:} الإعجاز ليس في النظم والبيان فحسب كما زعم بعض المتكلمين، بل في ما تضمنه من أخبار الغيب وأحكام التشريع والهدى (اللقاء الثالث).
    \item \textbf{الكلمات الأعجمية في القرآن:} ظاهرة وجود كلمات تستعملها العرب وتوجد في لغات أخرى (كالحبشية والفارسية)، وجهها الطبري بأنها من باب "اتفاق الألسنة"، فالعرب تنطق بها وغيرها ينطق بها بمعانيها، ولا يمنع ذلك من كون القرآن نزل بلسان عربي مبين (اللقاء الثالث).
    \item \textbf{معنى نزول القرآن على سبعة أحرف عند الطبري:} يرى الطبري أن الأحرف السبعة هي سبع لغات (ألسن) من لغات العرب، متفقة في المعنى ومختلفة في اللفظ (كالمترادفات: هلم، تعال، أقبل)، نزلت على وجه التيسير والتخيير، وليست هي القراءات السبع المشهورة (اللقاء الرابع).
    \item \textbf{الفرق بين "الأحرف السبعة" و"الأبواب السبعة":} يفرق الطبري بين حديث "سبعة أحرف" (ويرى أنها لغات)، وحديث "سبعة أبواب من الجنة" (ويرى أنها أوجه وموضوعات القرآن: كالأمر والنهي والحلال والحرام والمحكم والمتشابه والأمثال) (اللقاء الرابع).
    \item \textbf{سبب جمع القرآن على حرف واحد:} يرى الطبري أن الصحابة بقيادة عثمان بن عفان رضوان الله عليهم، لما رأوا الاختلاف والتنازع بين المسلمين في القراءة قد يفضي إلى التكفير، اختاروا جمع الأمة على حرف واحد (حرف قريش) درءاً للفتنة، وكان اختيارهم هذا من باب النظر للمصلحة العليا للإسلام وليس تضييعاً لما أُمروا بحفظه (اللقاء الرابع).
    \item \textbf{وجوه مطالب التأويل:} قسّم الطبري ما يُطلب تأويله إلى ثلاثة أقسام: ما لا يعلم إلا ببيان النبي ﷺ (تفصيل الشرائع)، وما استأثر الله بعلمه (كوقت الساعة)، وما يعلمه أهل لسان العرب (من إعراب ومعاني مفردات) (اللقاء الخامس).
    \item \textbf{الفرق بين الرحمن والرحيم:} "الرحمن" يدل على عموم الرحمة لجميع الخلق في الدنيا، بينما "الرحيم" يدل على خصوص الرحمة بالمؤمنين (اللقاء الخامس).
    \item \textbf{الترجيح بين القراءات المتواترة:} لا يكتفي الطبري بسرد القراءات بل قد يرجح إحداها لمعنى بلاغي أو سياقي، كترجيحه قراءة (مَلِك) على (مَالِك) لشموليتها ولئلا يتكرر معنى مالك المذكور سلفاً في كلمة (رب) (اللقاء السادس).
    \item \textbf{البسملة والفاتحة:} يرجح الطبري أن البسملة ليست آية من سورة الفاتحة، مستدلاً بعدم ورود تكرار آية متطابقة المعنى (الرحمن الرحيم) متجاورة بلا فاصل في القرآن (اللقاء السادس).
    \item \textbf{معنى العالمين:} المراد بالعالمين في سورة الفاتحة عند الطبري من خلال الآثار: الإنس والجن، أو كل صنف من الخلائق (اللقاء السادس).
    \item \textbf{تفسير الهداية بالتوفيق:} حصر الطبري الهداية المطلوبة في (اهدنا) بالتوفيق والثبات على الحق، وليس مجرد الإرشاد والبيان (اللقاء السادس).
    \item \textbf{جواز التسمية بسورة البقرة:} حديث "لا تقولوا سورة البقرة... ولكن السورة التي يذكر فيها البقرة" هو حديث منكر لا يثبت، وقد بوّب البخاري على جواز التسمية وأورد آثاراً صحيحة في ذلك (اللقاء السابع).
    \item \textbf{فواتح السور (الم):} مجرد اجتهاد الصحابة والتابعين في تفسير فواتح السور يدل على أنها من القرآن الذي يُعلم معناه وليست من المتشابه المطلق، ويرجح الطبري أنها حروف مقطعة تدل على معانٍ متعددة (كأسماء الله، والآلاء) وليست بلا معنى (اللقاء السابع).
    \item \textbf{الفرق بين (يَكْذِبُونَ) و (يُكَذِّبُونَ):} رجح الطبري قراءة تخفيف الذال (يَكْذِبُونَ) في وصف المنافقين؛ لأن السياق يعالج ادعاءهم الإيمان كذباً بألسنتهم، فالوعيد على هذا الكذب أليق بالسياق من الوعيد على مجرد التكذيب المطلق (اللقاء الثامن).
    \item \textbf{الترجيح بين القراءات بالمعنى التأسيسي:} يرجح الطبري القراءة التي تؤسس لمعنى جديد على القراءة التي تفيد التأكيد، كترجيحه (فأزَلَّهُما) بالتشديد بمعنى الاستذلال على (فأزَالَهُمَا) بمعنى التنحية والإخراج، لأن الله قال بعدها (فأخرجهما) (اللقاء العاشر).
    \item \textbf{قاعدة في رد القراءات لمخالفتها المستفيض:} قد يرد الطبري قراءة رغم صحتها من جهة لغة العرب (مثل قراءة نصب آدم ورفع كلمات في فتلقى آدم من ربه كلمات) إذا خالفت إجماع القرأة ونقل الحجة (اللقاء العاشر).
    \item \textbf{توجيه القراءات بالمعنى:} يرجح الطبري بين القراءات المتواترة أحياناً بناءً على السياق وقوة المعنى التأصيلي، ولا يكتفي بمجرد ثبوت السند (اللقاء الحادي عشر).
    \item \textbf{الاعتبار بقصص القرآن:} قصص السالفين في القرآن سيقت للعبرة ولتعليم الأمة، ليحذروا طريق الكفار ويتبعوا طريق الصادقين، وكتعزية للنبي ﷺ أن تكذيب اليهود له امتداد لعناد أسلافهم مع أنبيائهم (اللقاء الثاني عشر).
    \item \textbf{عدم توقف العبرة على تفاصيل الإسرائيليات:} القصص القرآنية تذكر مواضع العبرة، وما لم يُذكر من التفاصيل (كنوع الحجر الذي ضربه موسى) لا تتوقف عليه العبرة، فلا يُبنى عليها حكم ولا يضر جهلها (اللقاء الثاني عشر).
    \item \textbf{الانتفاع بقصص القرآن:} القصص القرآنية لم تُسق للتسلية ولا لمجرد التوثيق التاريخي، بل للاعتبار. وهذا الاعتبار لا يقع إلا من المتقين الذين طابت قلوبهم، فكل مسلم اعتبر بالقصص فهو المتقي المذكور في قوله: (وموعظة للمتقين) (اللقاء الثالث عشر).
    \item \textbf{تفسير الإيمان بالتصديق:} الإيمان يتضمن التصديق ولكنه أعلى منه، فهو يجمع بين الائتمان للمخبر والاتباع والتسليم لأمره، ولا يقتصر على التصديق اللغوي المجرد (اللقاء الرابع عشر).
    \item \textbf{بطلان دعوى الزيادة في القرآن:} لا يوجد في القرآن حرف أو كلمة زائدة لا تفيد معنى، وادعاء زيادة كلمات كـ (بأيديهم) أو (يطير بجناحيه) باطل مبني على توهم قياس نحوي قاصر، بل جاءت هذه الألفاظ لتأكيد المباشرة ونفي الاحتمال المجازي (اللقاء الرابع عشر).
    \item \textbf{الفرق بين العلم والمعرفة في القرآن:} المعرفة قد تكون مجرد الإدراك بلا عمل، أما العلم في السياق القرآني فيتضمن الإدراك المتبوع بالخشية والعمل والاتباع، لذا عرف اليهود الحق وكفرو به (اللقاء الخامس عشر).
    \item \textbf{حقيقة السحر وتعليم الملكين:} تعليم هاروت وماروت للسحر كان بامر الله ابتلاءً وامتحاناً للعباد، مع قيامهما بواجب التحذير، فالله يمتحن خلقه بما يشاء ويمحصهم ليميز الخبيث من الطيب (اللقاء الخامس عشر).
    \item \textbf{منهج الطبري في الترجيح وجمع المعاني:} من منهجه أنه يختار القراءة التي تستوعب دلالات المعاني المختلفة (كالإنساء والتأخير) متى أمكن ذلك، ويرد القراءات الشاذة المخالفة لرسم المصحف (اللقاء السادس عشر).
    \item \textbf{تخصيص اللفظ العام بالسياق والقرائن:} يميل الطبري إلى تخصيص اللفظ العام ببعض أفراده إذا وُجدت قرينة، كتخصيصه (الذين منعوا مساجد الله) بالنصارى، مستدلاً بالسياق السابق واللاحق الذي يذم أهل الكتاب (اللقاء السابع عشر).
    \item \textbf{الترجيح بين القراءات بالمعنى وقوة الحجة:} يرجح الطبري القراءة التي تتسق مع السياق ولا تعارضها الأحاديث الصحيحة، كترجيحه قراءة (وَلَا تُسْأَلُ) بالرفع على قراءة النهي لضعف سبب نزولها واتساقها مع السياق (اللقاء الثامن عشر).
    \item \textbf{منهج الطبري في تفاصيل الإسرائيليات:} يؤسس الطبري لعدم الجزم بتفاصيل القصص الغيبية (كقواعد البيت) التي لم تصح بها حجة، ويقرر أن "لا يضرنا الجهل بها كما لا ينفعنا العلم بها" ويتوقف فيها (اللقاء الثامن عشر).
    \item \textbf{الفرق بين أنواع الخلاف وأصنافه:} أنواع الخلاف ثلاثة (لفظي، تنوع، تضاد)، أما أصناف الخلاف فهي مجالاته (كالاختلاف في القائل، في عود الضمير، في سبب النزول، في الإعراب) (اللقاء التاسع عشر).
    \item \textbf{طريقة القرآن في الإخبار والتوجيه:} كل خبر في القرآن يتبعه توجيه للعمل والتصرف الصحيح، فالله يخبرنا بمكر الأعداء ليوجهنا للثبات والعمل الصحيح ولا نكتفي بمجرد المعرفة (اللقاء التاسع عشر).
\end{itemize}

\section{الفوائد الأصولية والمنهجية}
\begin{itemize}
    \item \textbf{الفرق بين العالم والإمام المحقق:} العالم هو كل من حصّل قدراً من العلم، أما الإمام المحقق فهو من كان له (أثر) في العلم تحريراً ونقداً، ومن أبرز علامات الإمام المحقق أنه "إذا فاتك كتابه يصعب جداً تعويضه" بغيره (اللقاء الأول).
    \item \textbf{أهمية دراسة سير الأئمة:} تفيد طالب العلم في معرفة كيفية نشأتهم، وتجعله يعرف قدر نفسه (فلا يغتر)، وتعلمه كيف يُميز بين الأصول والفضول، وتُعينه على فهم السياق الذي صُنف فيه الكتاب (اللقاء الأول).
    \item \textbf{خطورة الدخول في المطولات دون تأهب:} الكتب التأسيسية كالمطولات (تفسير الطبري، منهاج السنة، صحيح البخاري) لا تُقرأ جرداً، فالمقدمات التأهيلية ضرورية لفهم مصطلحات العالم، ولتعمل كـ "كشاف" لاصطياد الفوائد (اللقاء الأول).
    \item \textbf{مراعاة خاصة الكتاب (غايته):} لا يصح محاكمة كتاب لغاية لم يقصدها مؤلفه. الطبري قَصَد "بيان معاني القرآن" ولم يجعله موسوعة نحوية أو بلاغية إلا ما دعت إليه الحاجة، فلا يقال إن غيره أشمل منه في هذا الباب (اللقاء الثاني).
    \item \textbf{منهجية شرح وتدقيق التراث:} الشارح المحقق لا ينشغل بشرح وتفكيك كل لفظة في الكتاب (وإلا طال به المقام في الواضحات)، بل يتوقف عند الأمور التي تحتاج إلى فحص ودراسة وتعليق حقيقي (اللقاء الثاني).
    \item \textbf{بناء الخلاصات العلمية:} الخلاصات التي يكتبها العالم هي نتاج فهمه وجمعه لمن سبقه، ومعرفة مصدر العالم (كتأثر الطبري بالشافعي في الرسالة) يسهل جداً فهم مقصده (اللقاء الثاني).
    \item \textbf{كيفية قراءة مقدمات الكتب التأسيسية:} تعتمد على ثلاثة أمور رئيسية: حصر موضوعات المقدمة، حصر المسائل المندرجة تحتها، واستخراج الخلاصات والنتائج التي يريد المصنف الوصول إليها (اللقاء الثالث).
    \item \textbf{التفريق بين الظاهرة وتفسيرها:} يجب التمييز بين الظاهرة المتفق عليها (كوجود كلمات متوافقة بين اللغات) وبين التفسير وتوجيه هذه الظاهرة (كادعاء أخذ أحدهما من الآخر)؛ فالظاهرة حقيقة، أما تفسيرها فيحتاج إلى حجة وبرهان (اللقاء الثالث).
    \item \textbf{اختلاف التضاد والتنوع:} الخلاف الذي يمكن الجمع فيه باختلاف الجهة (كأن يُنسب لفظ للحبشية لكونهم ينطقون به، وللعربية لكونهم ينطقون به أيضاً دون نفي أحدهما للآخر) هو اختلاف تنوع، أما التضاد فهو ما لا يمكن فيه الجمع (اللقاء الثالث).
    \item \textbf{خطورة التسليم بالأصول الخاطئة:} من الخطأ المنهجي الجسيم التسليم للمخالفين بمقدمات وأصول باطلة (كإثبات أسلوب المشاكلة اللفظية أو المجاز المطلق) ثم محاولة الاستثناء منها في الفروع (كآيات الصفات)، بل الصواب هدم الأصول الباطلة من أساسها كما فعل شيخ الإسلام ابن تيمية (اللقاء الثالث).
    \item \textbf{منهج الطبري في عرض المسائل الكبرى:} من منهجه أنه يجمع كل ما ورد في الباب من روايات وآثار أولاً، ثم بعد استيفاء المادة العلمية يبدأ في تحليلها ونقدها وعرض رأيه فيها، ثم يستعرض الإشكالات الواردة على قوله ويجيب عنها (اللقاء الرابع).
    \item \textbf{سياق الحديث كقرينة للترجيح:} من أقوى حجج الطبري في تفسير حديث الأحرف السبعة هو سياق الحادثة نفسها (خصومة عمر وهشام بن حكيم)، حيث دل السياق على أن الخلاف كان في لفظ التلاوة لا في المعنى والحكم (اللقاء الرابع).
    \item \textbf{توجيه آثار السلف في الإحجام عن التفسير:} لم يحجم السلف (كسعيد بن المسيب وعبيدة السلماني) عن التفسير لكونه محرماً، بل تورعاً وخوفاً من ألا يبلغوا باجتهادهم إصابة الصواب، قياساً على إحجامهم عن الفتيا (اللقاء الخامس).
    \item \textbf{آلة المفسر وحدود اجتهاده:} يشترط الطبري في المفسر العلم بالآثار النبوية الصحيحة، والبراهين اللغوية من الشواهد، وألا يخرج في تأويله عن مجموع ما قالته الحجة من طبقات الصحابة والتابعين (اللقاء الخامس).
    \item \textbf{تأخير المسائل لمظانها:} من منهج الطبري أنه لا يتوسع في مسألة أو يفرّع فيها في المقدمات إذا كان لها موضع أليق وأنسب في ثنايا التفسير العملي للآيات، بل يعزوها إلى مكانها (اللقاء الخامس).
    \item \textbf{الإثم في التفسير بالرأي بلا علم:} القائل في القرآن بغير علم ولا حجة، حتى وإن وافق قوله الصواب صدفة، فهو مخطئ وآثم لتقوله على الله بغير برهان (اللقاء الخامس).
    \item \textbf{تفسير الطبري ليس بالمأثور فقط:} يخطئ من يحصر تفسير الطبري في الجمع والمأثور؛ بل هو إمام مجتهد يورد الأقوال، وينقدها، ويطرح الإشكالات ويجيب عنها، ويوجه الإعراب لخدمة المعنى (اللقاء السادس).
    \item \textbf{منهجية الطبري في تأخير النقض:} من طريقته أنه قد يورد قولاً مخالفاً ويسرد حجته كاملة فلا ينبغي التعجل بظن أنه يرتضيه، فقد يؤخر الرد عليه وتخطئته صفحات عديدة (اللقاء السادس).
    \item \textbf{الدمج بين الشواهد اللغوية والآثار:} يبدأ الطبري في تفسير المفردة (ككلمة رب) بلسان العرب والشعر، ثم يثني بآثار الصحابة، ليجمع بين التفسير اللغوي والأثري (اللقاء السادس).
    \item \textbf{إطراح الروايات شديدة الضعف:} الطبري قد يروي الضعيف للاستئناس اللغوي، ولكنه يطرح الرواية تماماً ولا يعتمدها (كرواية حساب الجُمَّل في فواتح السور) إذا تيقن بطلانها وتهالك رواتها (اللقاء السابع).
    \item \textbf{الإلزام في المناظرة:} من طرائق الطبري في الرد على المخالفين إلزامهم بأصول يقرون بها؛ كإلزامه لمن أنكر دلالة الحرف الواحد على معانٍ متعددة بأنهم يقرون بدلالة الكلمة الواحدة على معانٍ متعددة (اللقاء السابع).
    \item \textbf{بيان لوازم الأقوال والتوجيهات:} يبرز الطبري النتائج المترتبة على الأقوال التفسيرية لبيان قوتها أو ضعفها (كالقول بأن "ذلك" إشارة للتوراة، مما يدفع إشكال مجيء "ذلك" للغائب بدلاً من "هذا" للحاضر) (اللقاء السابع).
    \item \textbf{اختلاف التضاد والتنوع:} التأسيس للتفريق بين اختلاف الأقوال الذي يمكن الجمع فيه لاختلاف الجهة (تنوع)، والاختلاف الذي يستحيل فيه الجمع (تضاد) (اللقاء السابع).
    \item \textbf{تأثر الطبري بمدرسة النحاة في التقدير:} من المآخذ المنهجية تأثر الطبري ببعض النحويين (كالفراء وأبي عبيدة) في افتراض قياس نحوي معين للنص القرآني، وجعل ما خالفه من باب الحذف أو الزيادة، مما استغله المتكلمون لاحقاً في تأويلات فاسدة (اللقاء الثامن).
    \item \textbf{الفرق بين النزول والتنزّل:} يفرق المحققون بين "نزول" الآية (فيمن نزلت تاريخياً) وبين "تنزلها" وانسحاب حكمها على من شاركهم في الوصف عبر الزمان (اللقاء الثامن).
    \item \textbf{التثبت من مأخذ القول قبل نقده:} من الخطأ بناء نقد على افتراض نية أو علة للعالم لم يصرح بها، كافتراض الطبري أن مجاهداً خص آية بأهل الكتاب لجهله بعلم العرب بربوبية الله، وهو إمام جليل لا يجهل ذلك (اللقاء التاسع).
    \item \textbf{موقف الطبري وابن كثير من الإسرائيليات:} الطبري يميل إلى إيراد الروايات الغريبة والإسرائيلية دون تعقيب ما لم تخالف ظاهراً، بينما ابن كثير يميل إلى نقدها وتضعيفها ووصفها بالنكارة (اللقاء التاسع).
    \item \textbf{الاكتفاء بصحيح الوحي في الوعظ:} لا يحتاج الداعية لترغيب الناس وترهيبهم إلى الأحاديث الضعيفة والموضوعة والإسرائيليات الغريبة، ففي محكم القرآن وصحيح السنة الغنى والكفاية (اللقاء التاسع).
    \item \textbf{تحرير محل الخلاف:} عند اختلاف المفسرين في أمر معلوم من الدين بالضرورة (كعلم الله بما لا نعلم)، فإنهم لا يختلفون في الأصل، بل يختلفون في تحديد "الغائب المعين" الذي يشير إليه السياق (اللقاء العاشر).
    \item \textbf{قاعدة في نقد الأقوال:} العالم المحقق إذا أراد رد قول فإنه يوجهه أولاً مبيناً حجته، ثم ينقضه ببيان لوازمه الباطلة ومخالفته للغة وإجماع المفسرين (اللقاء العاشر).
    \item \textbf{منهجه في المبهمات التي لا طائل منها:} يؤسس الطبري لعدم الانشغال بتفاصيل المبهمات (كنوع الشجرة التي أكل منها آدم)، لأن العلم بها لا ينفع والجهل بها لا يضر، ولو كان في معرفتها خير لبينه الله (اللقاء العاشر).
    \item \textbf{منهجه في نقد الإسرائيليات:} يطرح الطبري تفاصيل الروايات الإسرائيلية إذا خالفت ظاهر القرآن (كقصة دخول إبليس في فم الحية لمعارضة قوله تعالى "وقاسمهما")، مقرراً عدم ترك المفهوم الظاهر إلى باطل لا دلالة عليه (اللقاء العاشر).
    \item \textbf{توجيه الآثار المتعارضة ظاهرياً:} يجمع الطبري بين الأقوال إن أمكن الجمع باختلاف الجهة، كاعتباره إبليس من الملائكة (أنه من حي منهم يُقال لهم الجن خُلقوا من نار) ليجمع بين استثنائه من الملائكة وكونه من الجن (اللقاء العاشر).
    \item \textbf{تحرير موضع الإجماع والاختلاف:} عند ذكر مسألة، يحرص الطبري على بيان القدر المتفق عليه بين العلماء، ثم تحديد النقطة الدقيقة التي وقع فيها الخلاف لتوجيه النقاش (اللقاء الحادي عشر).
    \item \textbf{صحة المعنى لا تستلزم صحة التفسير:} قد يكون المعنى المذكور صحيحاً في ذاته (كأن يقال إن الكفر بالقرآن هو كفر بمحمد وبالتوراة)، لكن هذا لا يلزم أن يكون هو التفسير المراد من الآية المحددة في سياقها (اللقاء الحادي عشر).
    \item \textbf{تأخير الردود والمناقشات:} قد يذكر الطبري قوله أو قول غيره في مسألة مجْمَلاً مع أدلته، ويؤخر تفصيل الحجاج والرد على المخالفين إلى موضع أليق في القرآن أو بعد صفحات (اللقاء الحادي عشر).
    \item \textbf{تخصيص دلالة الألفاظ العامة بالسياق:} الألفاظ العامة كـ (التقوى، البر، الإثم) لا يبقيها الطبري على عمومها المطلق دائماً حين التفسير، بل يفسرها بما يلائم السياق الخاص الذي وردت فيه (اللقاء الحادي عشر).
    \item \textbf{تقديم ما قرب ذكره:} من قرائن الترجيح عند الطبري تقديم عودة الضمير أو الصفة على ما قرب ذكره في السياق دون ما بَعُد، ما لم توجد قرينة تصرفه (اللقاء الحادي عشر).
    \item \textbf{عدم القطع في موضع الظن:} الطبري لا يقطع بصحة تنزيل حديث نبوي عام على آية معينة (كحديث الطاعون وتنزيله على الرجز) ما لم يكن هناك دليل قاطع أو تصريح بذلك، بل يبقيه في دائرة الاحتمال (اللقاء الحادي عشر).
    \item \textbf{موقف الطبري عند تكافؤ الأدلة (عدم القطع في الظن):} إذا اختلفت الأقوال (مثل هل "مصراً" هي مصر المعروفة أم أي مِصر من الأمصار) ولم يوجد خبر يقطع العذر، فإن الطبري لا يجزم في موضع الظن، بل يكتفي بالقدر المحكم ويعمم ما عممه القرآن (اللقاء الثاني عشر).
    \item \textbf{رد الإسناد لنكارة المتن:} قد يورد الطبري إسناداً ويرده لكون متنه منكراً أو غريباً وليس له شاهد، حتى وإن بدا الإسناد متصلاً في ظاهره (اللقاء الثاني عشر).
    \item \textbf{منهجية تأخير الترجيح وشرح الاحتمالات:} من طريقة الطبري أنه يورد احتمالات الكلمة (كعودة الضمير)، ثم يشرح معنى الآية بناءً على كل احتمال، ثم يتأخر جداً في ذكر اختياره وترجيحه وتخطئة الباقي (اللقاء الثالث عشر).
    \item \textbf{قاعدة في رد الأقوال الشاذة:} لا يجوز ترك المفهوم من ظاهر التنزيل إلى معنى باطن (كتأويل مجاهد لمسخ القردة بأنه مسخ قلوب) بغير دلالة من كتاب ولا سنة صحيحة ولا إجماع، ومخالفة إجماع المفسرين تكفي كدليل على فساد القول (اللقاء الثالث عشر).
    \item \textbf{مهارة التلخيص في دراسة المطولات:} التلخيص دليل الفهم، وهو القدرة على بناء خرائط ذهنية واستخراج الجمل المركزية والنتائج من كلام المصنف وحججه للترجيح، بدلاً من التيه في الروايات (اللقاء الرابع عشر).
    \item \textbf{تخصيص اللفظ العام بالقرينة:} الطبري قد يخصص اللفظ العام ببعض أفراده إذا دلت عليه قرينة سياقية، كتخصيصه "السيئة" بالشرك لذكر الخلود في النار معها، حيث أن الخلود لا يقع بغير الشرك (اللقاء الرابع عشر).
    \item \textbf{أهمية السياق الزمني والمكاني:} قراءة الحدث في سياقه الزمني يزيل الإشكالات ويبين وجه قوة الحجة، كإخبار النبي ﷺ بأسرار اليهود وهم سادة العلم في ذلك الزمان (اللقاء الخامس عشر).
    \item \textbf{معرفة مصادر الأقوال:} معرفة المصادر اللغوية التي استقى منها الطبري (كأبي عبيدة والفراء والأخفش) تُعين على فهم وجه الإيرادات النحوية واللغوية في تفسيره وكيفية تعامله معها (اللقاء الخامس عشر).
    \item \textbf{العفو والصفح بين النسخ والإحكام:} ذهب بعض أهل العلم إلى أن آيات العفو والصفح نُسخت بآيات السيف. والراجح أنها آيات محكمة يُعمل بها بحسب الحال من الاستضعاف والقوة، والنسخ عند المتقدمين أوسع من إزالة الحكم بالكلية (اللقاء السابع عشر).
    \item \textbf{الجمع بين الأحاديث المتعارضة ظاهرياً:} لا يجوز في أخبار النبي ﷺ الثابتة أن يدفع بعضها بعضاً، بل يجب الجمع بينها؛ كجمع الطبري بين حديث تحريم الله لمكة وحديث تحريم إبراهيم لها باختلاف الجهة (اللقاء الثامن عشر).
    \item \textbf{منهج الطبري في عدم الترجيح بلا دليل:} لا يجزم بتعيين المبهمات طالما لم يرد نص قاطع، وينقد المراسيل ويرد الأحاديث التي في أسانيدها نظر إذا خالفت السياق (اللقاء الثامن عشر).
    \item \textbf{تخطئة الأقوال اللغوية لمخالفتها الإجماع:} يرد الطبري تأويلات النحويين (كالفراء والأخفش وأبي عبيدة) بشدة إذا خالفت إجماع المفسرين من الصحابة والتابعين وسياق الآيات (اللقاء التاسع عشر).
\end{itemize}

\section{الفوائد اللغوية}
\begin{itemize}
    \item \textbf{مركزية علوم العربية:} مدار دراسة فهم آيات القرآن يرتكز على العناية بلغة العرب. ومن جهلها دخلت عليه الشبهات واللبس، وهو المعنى الذي أكده الإمام الشافعي وتأثر به الطبري في مقدمته (اللقاء الثاني).
    \item \textbf{أهمية الاعتماد على نسخ مضبوطة بالشكل:} كلام الأئمة دقيق ومكثف، ولا يُفهم على وجهه إلا بضبط الكلمة، كالتفريق بين المبني للمعلوم والمبني للمجهول (مثل قوله: وخُدِعَ الشيطانُ) لتحديد الفاعل من المفعول (اللقاء الثاني).
    \item \textbf{عدم طرد الأساليب العربية بلا قرينة:} أساليب العرب في الخطاب (كالمشاكلة اللفظية) لا يصح طردها وتطبيقها دائماً في القرآن إذا خالفت مقدمات عقدية قطعية أو سياقاً محدداً (اللقاء الثالث).
    \item \textbf{إطلاق المصدر على المفعول:} سُمي القرآن قرآناً (وهو مصدر بمعنى القراءة) ويُقصد به المقروء، كقولهم "الكتاب" ويُقصد به "المكتوب" (اللقاء الخامس).
    \item \textbf{قاعدة في ترتيب الأسماء والصفات:} من شأن العرب إذا أخبرت عن شيء أن تُقدم الاسم العَلَم أولاً، ثم تُتبعه بالصفات والنعوت، ولذلك قُدم اسم الجلالة "الله" على "الرحمن الرحيم" (اللقاء الخامس).
    \item \textbf{الاحتجاج بالضعيف لغوياً:} قد يورد الطبري أثراً ضعيف الإسناد لا ليصححه حديثياً، بل للاستشهاد به على صحة معنى الكلمة في لغة العرب (اللقاء السادس).
    \item \textbf{دلالة الصلاة لغوياً:} تأثر بعض اللغويين بأن الصلاة لغة هي الدعاء الخالص، والصواب أنها في كلام العرب أعم من ذلك وهي عبادة ذات هيئة تشمل الدعاء والتذلل (اللقاء السابع).
    \item \textbf{دلالة صيغة المفاعلة:} رد الطبري على من زعم أن "يخادعون" تأتي من طرف واحد، مبيناً أن صيغة التفاعل تقتضي المشاركة، فالمنافق يخادع، والله يخادعه بالاستدراج والمجازاة (اللقاء السابع).
    \item \textbf{الرد على تكلف اللغويين:} استنكار اللغويين لمجيء حرف بدلاً من آخر (مثل: خلوا إلى شياطينهم بدلاً من بشياطينهم) هو تكلف؛ فالنص القرآني جارٍ على سعة كلام العرب، وافتراض "أصل" يجب القياس عليه دون غيره منشأ لكثير من الشبهات (اللقاء الثامن).
    \item \textbf{خطورة جعل القواعد النحوية حاكمة على القرآن:} القرآن نزل بلسان عربي مبين وهو الميزان، فلا يصح بناء قاعدة لغوية ضيقة ثم اتهام النص القرآني بالمجاز أو الحذف والزيادة لكي يوافقها (اللقاء التاسع).
    \item \textbf{دلالة لفظ التسبيح:} لا يقتصر التسبيح في كلام العرب على التنزيه ونفي النقص، بل يشمل التعظيم والذكر والصلاة (اللقاء العاشر).
    \item \textbf{تغليب الأغلبية (التغليب):} استدل الطبري بلفظ (عَرَضَهُم) على أن المسميات هي للملائكة وبني آدم، لأن العرب تكنّي بالميم والهاء عن العقلاء (اللقاء العاشر).
    \item \textbf{المنع من الصرف للعُجمة:} كلمة (إبليس) فِعْليل من الإبلاس، ومُنعت من الصرف استثقالاً لشبهها بالأسماء الأعجمية التي لا نظير لها في كلام العرب (اللقاء العاشر).
    \item \textbf{القرآن ميزان اللغة وليس العكس:} القرآن نزل بلسان عربي مبين، وهو الحجة في إثبات عربية الألفاظ والأساليب، فمن الخطأ افتراض قياس نحوي ثم اتهام نص القرآن بالحذف أو الزيادة لموافقته (اللقاء الحادي عشر).
    \item \textbf{أصل تسمية "النبي":} يرى الطبري أن "النبي" أصلها مهموز من "النبأ" (لأنه مُنبأ من الله ومُنبئ عنه)، وصُرفت إلى "فعيل"، وقيل من "النبوة" أي الارتفاع، والقول الأول أقوى (اللقاء الثاني عشر).
    \item \textbf{معنى القتل (بغير الحق):} قول القرآن (ويقتلون النبيين بغير الحق) ليس معناه إقرار أن هناك قتلاً للأنبياء "بحق"، بل هو بيان لحالهم ومبالغة في وصف تجردهم عن أي حجة في اعتدائهم (اللقاء الثاني عشر).
    \item \textbf{الفرق بين العلم والمعرفة:} المعرفة أقل من العلم؛ فهي إدراك مجرد قد لا يتبعه عمل، أما العلم في السياق القرآني والشرعي فيتضمن الإدراك المتبوع بالخشية والعمل والاتباع (اللقاء الثالث عشر).
    \item \textbf{معنى الخاسئ:} الخاسئ هو المبعد المطرود الذليل الصاغر، كما يُخسأ الكلب (اللقاء الثالث عشر).
    \item \textbf{أصل النكال:} النكال مصدر من التنكيل، وأصله العقوبة (اللقاء الثالث عشر).
    \item \textbf{معنى الأميين:} الراجح أن الأميين هم من ليس لهم كتاب (كالعرب)، تمييزاً لهم عن أهل الكتاب، وليس مجرد من لا يقرأ ولا يكتب (اللقاء الرابع عشر).
    \item \textbf{الاستثناء المنقطع:} مجيء أداة الاستثناء بمعنى "لكن"، كقوله (إلا أماني) أي لكنهم يختلقون الأكاذيب، لأن الأماني ليست من جنس الكتاب المذكور قبلها (اللقاء الرابع عشر).
    \item \textbf{أصل الفتح:} الفتح في كلام العرب يأتي بمعنى النصر، ويأتي بمعنى القضاء والحكم كما في قوله (بما فتح الله عليكم) أي بما قضى وحكم به (اللقاء الرابع عشر).
    \item \textbf{عدم وجود حروف زائدة للصلة:} يرفض الطبري دعوى اللغويين بوجود حروف "صلة" زائدة لا معنى لها في القرآن، ويؤكد أن كل حرف له دلالة مقصودة (اللقاء الخامس عشر).
    \item \textbf{الخلاف النحوي لا يبطل المعنى المستقر:} لا يصح إبطال المعنى المستقر من ظاهر التنزيل بناءً على قياس نحوي مجرد أو ادعاء تناوب الحروف بلا حجة (اللقاء الخامس عشر).
    \item \textbf{إطلاق لفظ الزوج على المرأة:} الأفصح في لغة الحجاز والقرآن استعمال (الزوج) للذكر والأنثى بلا تاء تأنيث، و(زوجة) لغة تميم (اللقاء السادس عشر).
    \item \textbf{تضمين الأفعال معنى القسم:} قد يُشرب الفعل معنى اليمين فيُعامل معاملته نحوياً، كفعل (علموا) في قوله (ولقد علموا لمن اشتراه)، فجاءت اللام واقعة في جواب القسم المقدر (اللقاء السادس عشر).
    \item \textbf{دلالة (أم) المنقطعة:} تأتي (أم) للاستفهام المنقطع (المستأنف) بمعنى (بل)، فتفيد التوبيخ والإنكار ولا تفيد الشك أو التسوية (اللقاء السادس عشر).
    \item \textbf{بطلان الزيادة في القرآن بلا معنى:} لا يوافق المحققون كالطبري وابن تيمية على أن يكون في القرآن حرف "زائد" لا معنى له (وجوده وعدمه سواء). وما أُطلق عليه "زائد" مقصوده أنه يمكن استقامة الجملة إعرابياً دونه، لكن وجوده يضيف معنى وتأكيداً (اللقاء السابع عشر).
    \item \textbf{إطلاق الحجة على الدعوى الباطلة:} لفظ "حجة" قد يُطلق على الخصومة والدعوى ولو كانت باطلة (لئلا يكون للناس عليكم حجة إلا الذين ظلموا)، فليس كل حجة صحيحة (اللقاء التاسع عشر).
\end{itemize}

\section{الفوائد العقدية والإيمانية}
\begin{itemize}
    \item \textbf{الغاية من التفقه والمدارسة:} أعظم ما تطلبه النفس وتنشغل به في الدنيا هو العلم بربها، والعلم بوحي الله، وما خُلقت له، والاستقامة على ذلك (اللقاء الأول).
    \item \textbf{حفظ الدين ونصرته:} بث العلم وتعليمه للأمة لا يعادله شيء من الأعمال بعد النبوة، وهو من أعظم صور نصرة الله ورسوله صلى الله عليه وسلم (اللقاء الأول).
    \item \textit{(تضاف الفوائد العقدية المتعلقة بالأسماء والصفات ونقد الفرق عند دراسة الآيات إن شاء الله).}
    \item \textbf{نقد مصطلح "المعجزة":} لم يرد هذا اللفظ في القرآن ولا في السنة، وإنما الوارد هو (الآيات، البينات، البصائر، السلطان، البرهان). والآيات غايتها الهداية وبيان الحق وإقامة الحجة، وليست مجرد "التعجيز" و"التحدي" الابتدائي كما قرره المتكلمون (اللقاء الثالث).
    \item \textbf{بطلان القول بـ "الصَّرْفَة":} قول المعتزلة بأن العرب كانوا قادرين على الإتيان بمثل القرآن ولكن الله صرفهم بقوته هو قول باطل، وفيه غلو في قدرة الإنسان، وهو مأخوذ من فلاسفة اليونان (اللقاء الثالث).
    \item \textbf{المراء في القرآن كفر:} ورد في الحديث أن الجدال والمراء في القرآن (أي في ثبوته وألفاظه المنزلة) كفر، وهذا يدل على خطورة التشكيك في أصل الوحي المنزل (اللقاء الرابع).
    \item \textbf{إثبات خلق أفعال العباد (الرد على القدرية):} استنبط الطبري من (وإياك نستعين) افتقار العبد لعون الله، ومن (الضالين) أن إضافة الفعل للعبد لا تنفي أن الله خلقه وأوجده وأضله، وهذا يهدم أصول القدرية القائلين باستقلال العبد بفعله (اللقاء السادس).
    \item \textbf{إثبات صفات الأفعال لله تعالى:} أثبت الطبري صفة "الغضب" لله كما تُفهم في اللغة بلا تأويل (كردّه على من أولها بالعقوبة)، وبلا تمثيل (حيث نفى عنها الانفعال والاضطراب البشري)، قياساً على صفتي العلم والقدرة اللتين يثبتهما المخالف (اللقاء السادس).
    \item \textbf{الختم على القلوب عقوبة إلهية:} رد الطبري بقوة على المعتزلة الذين أوّلوا "ختم الله" بأنه "تكبر الكفار"، مبيناً أن الختم فعل لله كعقوبة مستحقة على توالي ذنوبهم (اللقاء السابع).
    \item \textbf{تحرير مسألة تكليف ما لا يطاق:} استدل الطبري بآية الختم على جواز تكليف ما لا يطاق، والصواب أن الله لا يكلف نفساً إلا ما تقدر عليه ابتداءً، أما سبق علم الله بأنهم لن يؤمنوا وعقوبته لهم بالختم فلا يُسمى تكليفاً لما لا يُطاق من جهة القدرة (اللقاء السابع).
    \item \textbf{الإيمان قول وعمل:} الإيمان شرعاً كلمة جامعة للإقرار والتصديق القلبي، وتصديق ذلك بالعمل، ورد بذلك على غلاة المرجئة الذين يحصرونه في التصديق اللغوي المجرد (اللقاء السابع).
    \item \textbf{مفهوم العذر بالجهل:} لا يُعاقب العبد إلا بعد أن تبلغه حجة الله، ومعنى (يَحْسَبُونَ أَنَّهُم مُّهْتَدُونَ) ليس أنهم كانوا يجهلون الحق، بل أنهم توهموا أن طريق النفاق الذي اختاروه يحقق لهم مصالحهم وينجيهم (اللقاء الثامن).
    \item \textbf{إثبات صفات الأفعال لله كالاستهزاء:} أثبت الطبري صفة "الاستهزاء والمكر" لله حقيقةً على وجه الكمال والمجازاة، ورد بقوة على من أوّلها باللعب والعبث أو حصرها في كونها مجرد "عقوبة" (اللقاء الثامن).
    \item \textbf{أسباب الإضلال والرد على الجبرية:} إضلال الله للعبد يكون حقاً وعدلاً بسبب من العبد وكسبه، كما نص القرآن (وما يضل به إلا الفاسقين)، وفي ذلك رد صريح على الجبرية (اللقاء التاسع).
    \item \textbf{إثبات صفة العلو ومنهج الإلزام:} أثبت الطبري استواء الله وعلوه حقيقة بلا تعطيل، وألزم المعطلة بأن فرارهم من الإثبات بحجة التنزيه سيوقعهم في تشبيه وتحريف أشنع في الألفاظ التي أولوها إليها (اللقاء التاسع).
    \item \textbf{بطلان مقولة "آدم خليفة الله":} لا يصح إطلاق هذا الوصف المبتدع، فالله هو الغني الشهيد ولا يغيب ليخلفه أحد، وإنما آدم وذريته يخلف بعضهم بعضاً في الأرض، أو يخلفون من سبقهم (اللقاء التاسع).
    \item \textbf{عقوبات الله بأسباب العباد:} إذا ذكر الله عقاباً في الدنيا أو الآخرة أرجعه إلى فعل الكفار أنفسهم وتعديهم (ذلك بما عصوا)، رداً على الجبرية، فالله لا يظلم أحداً مثقال ذرة (اللقاء الثاني عشر).
    \item \textbf{استعمال النعم في معصية الله:} من أشد الجحود أن يستعمل العبد ما وهبه الله من نعم في معصية المنعِم، وهذا من صور الكفر العملي بنعمة الله (اللقاء الثاني عشر).
    \item \textbf{الإيمان قول وعمل:} الإيمان في الشرع هو اعتقاد وقول يصدقه الاتباع والعمل، ولا يكتفى فيه بمجرد التصديق اللغوي المجرد، وفي ذلك رد على المرجئة (اللقاء الثاني عشر).
    \item \textbf{الابتلاء بقرب المعصية وسهولتها:} من سنن الله ابتلاء عباده بتيسير أسباب المعصية وقربها من أيديهم ليعلم من يخافه بالغيب، كما ابتلى أصحاب السبت بظهور الحيتان (اللقاء الثالث عشر).
    \item \textbf{دلالة نفي النقص عن الله:} كل نفي لنقص عن الله تبارك وتعالى (كقوله: وما الله بغافل عما تعملون) يتضمن إثبات كمال الضد (وهو كمال العلم والإحاطة والمراقبة) (اللقاء الثالث عشر).
    \item \textbf{العلم بالله مفتاح الهداية والطمأنينة:} كلما علم العبد عن ربه وأسمائه وصفاته ازداد هداية ويقيناً وطمأنينة، وعلم أنه محيط بالكافرين وليس بغافل عنهم (اللقاء الرابع عشر).
    \item \textbf{عقوبة الإعراض عن الآيات:} من عقوبات العناد والإعراض عن الحق بعد تبينه أن يُصرف العبد عن الإيمان ويُختم على قلبه، فلا تنفعه الآيات بعدها (اللقاء الرابع عشر).
    \item \textbf{الرد على الخوارج والمعتزلة في الشفاعة:} الآيات التي تذكر الخلود في النار لفاعلي السيئات مخصوصة بأهل الشرك والكفر، ولا تدخل فيها الكبائر التي دون الشرك لثبوت الشفاعة لأهلها خلافاً لمنكري الشفاعة (اللقاء الرابع عشر).
    \item \textbf{أمانة العلم:} العلم ميثاق ومسؤولية من الله، ومن استعمله للبغي أو الحسد أو نيل الدنيا فقد شابه اليهود الذين اشتروا الضلالة بالهدى (اللقاء الخامس عشر).
    \item \textbf{دقة الحجاج القرآني:} إفحام المخالفين بالنقض من كلامهم؛ كاحتجاجه على اليهود بقتلهم الأنبياء وهم يدعون الإيمان بالكتاب الذي يوجب اتباعهم (اللقاء الخامس عشر).
    \item \textbf{التفريق بين التمحيص والإهانة:} العذاب المهين مختص بالكفار، أما ما يصيب عصاة المؤمنين فهو تمحيص لذنوبهم، فالله يطهر وليه ولا يهينه (اللقاء الخامس عشر).
    \item \textbf{الفرق بين الإذن الكوني والإذن الشرعي:} الإذن الشرعي هو ما يحبه الله ويشرعه، والإذن الكوني (القدري) هو ما يقدره الله بحكمته وإن لم يحبه، كالإذن بضرر السحر (اللقاء السادس عشر).
    \item \textbf{تهذيب الألفاظ وسد الذرائع:} نهي الله المؤمنين عن قول (راعنا) لاحتمالها معنى فاسداً عند اليهود أصل في سد الذرائع وتجنب الألفاظ الموهمة توقيراً للأنبياء (اللقاء السادس عشر).
    \item \textbf{الحسد من عند أنفسهم وتبرئة الوحي:} في قوله (حسدا من عند أنفسهم) تبرئة للتوراة من أن تكون أمرت بذلك، بل هو حسد من أنفسهم، وهذا أصل في التفريق بين الدين الصافي وأفعال المنتسبين إليه، وفيه أن أبواب الباطل لا تقتصر على الجهل بل تشمل مرض القلب (اللقاء السابع عشر).
    \item \textbf{عقوبة رد الحق بعد تبينه:} كفرهم لم يكن عن جهل بل عناد (من بعد ما تبين لهم الحق)، وعقوبة من يعاند الحق أن يُصرف عنه ويُعمى قلبه (اللقاء السابع عشر).
    \item \textbf{الرد على المعتزلة في قدرة الله:} قوله (إن الله على كل شيء قدير) عموم لا يدخله التخصيص، وفيه رد على المعتزلة الذين زعموا أن الله لا يقدر على أفعال العباد (اللقاء السابع عشر).
    \item \textbf{أمانة العلم وتلاوة الكتاب:} وبخ الله اليهود والنصارى لأنهم يقرأون الكتاب المصدق للأنبياء وسوى بينهم وبين الجاهلين، مما يوجب على طالب العلم أن يكون وقافاً عند كتاب الله وإلا كان والجاهل سواء (اللقاء السابع عشر).
    \item \textbf{الرد على المتكلمين في (كن فيكون):} استشكل المتكلمون أمر المعدوم بـ "كن"، والصواب أن أمر الله للشيء بالوجود لا يتقدمه ولا يتأخره، والمشكلة في افتراضهم قواعد عقلية قاصرة لمحاكمة النصوص بينما السلف لم يستشكلوا ذلك لعلمهم بكمال قدرة الله (اللقاء السابع عشر).
    \item \textbf{الاطلاع على بواطن الكفار يقطع التعلق بهم:} إعلام الله نبيه بأن أهل الكتاب لن يرضوا عنه يقطع الطمع في إرضائهم، ويوجه الداعية للتمسك بهدى الله وعدم المداهنة (اللقاء الثامن عشر).
    \item \textbf{أمانة العلم وعقوبة كتمانه:} من أُعطي العلم واتبع هواه (كأحبار اليهود) أو كتم الحق طلباً للرئاسة فحسابه أعظم، وهو متوعد بانقطاع ولاية الله ونصرته (اللقاء الثامن عشر).
    \item \textbf{معنى التلاوة الحقة:} (يتلونه حق تلاوته) تعني اتباعه والعمل به وتحليل حلاله وتحريم حرامه، لا مجرد تجويد ألفاظه كما اشتهر (اللقاء الثامن عشر).
    \item \textbf{الإيمان تصديق واتباع:} لا يكون التصديق إيماناً شرعياً نافعاً إلا بالاتباع، فمجرد المعرفة بالكتاب لا تغني بدون العمل (اللقاء الثامن عشر).
    \item \textbf{عصمة الأنبياء من الذنوب:} إجماع دعاء الأنبياء بتوبة الله عليهم يدل على احتياجهم للتوبة في الصغائر أو الأخطاء (كذنب آدم)، وهذا لا يقدح في النبوة بل الله يسددهم ويرفعهم بالتوبة الصادقة (اللقاء الثامن عشر).
    \item \textbf{أفضلية هدي الأنبياء:} من ادعى أن طريقاً غير طريق الرسل أهدى سبيلاً فقد ضل، فخير الناس هم الرسل، وخير الهدى ما جاؤوا به (اللقاء الثامن عشر).
    \item \textbf{التسليم المطلق لأحكام الله:} العبد يسلم لأوامر الله سواء علم الحكمة منها أم لم يعلم، وربط الإيمان بمعرفة الحكمة هو عبادة لله على حرف (اللقاء التاسع عشر).
    \item \textbf{الابتلاء بالتشريع:} كما يبتلي الله عباده بالمصائب والمنع والمنح، يبتليهم بالتشريع وتغيير الأحكام (كتحويل القبلة) ليميز من يسلم ممن يعترض (اللقاء التاسع عشر).
    \item \textbf{إثبات العلم الحادث لله (لنعلم):} العلم الحادث (لنعلم واقعاً) لا ينفي سبق العلم الأزلي المقدر، فالله يحاسب على ما وقع. ولا حاجة لتأويل الطبري بصرف الضمير لغير الله فراراً من إشكالات المتكلمين (اللقاء التاسع عشر).
    \item \textbf{توجيه أوامر القرآن للنبي ﷺ:} الأوامر والنواهي في القرآن للنبي ﷺ (كقوله: فلا تكونن من الممترين) موجهة إليه أصالة، وتعليق الحكم على شرط لا يستلزم وقوعه، وتأويل الطبري بأن الخطاب مراد به أمته محل نقد (اللقاء التاسع عشر).
    \item \textbf{الشكر لا يكون تاماً إلا بالعمل:} فسر الطبري الشكر بالثناء، وفيه قصور؛ فالشكر شرعاً لا يكتمل إلا بالعمل واستعمال النعم في طاعة المنعِم (اللقاء التاسع عشر).
\end{itemize}

\section{الأحكام الشرعية والفقهية}
\begin{itemize}
    \item \textit{(تضاف الفوائد الفقهية واستنباطات الأحكام الشرعية عند تناول آيات الأحكام إن شاء الله).}
\end{itemize}